\documentclass[11pt]{amsart}
\usepackage{amsmath,amssymb,amsthm}
\usepackage[margin=1.1in]{geometry}
\usepackage{booktabs}
\usepackage{hyperref}

\theoremstyle{plain}
\newtheorem{theorem}{Theorem}[section]
\newtheorem{proposition}[theorem]{Proposition}
\newtheorem{corollary}[theorem]{Corollary}
\newtheorem{lemma}[theorem]{Lemma}
\newtheorem{conjecture}[theorem]{Conjecture}
\theoremstyle{definition}
\newtheorem{definition}[theorem]{Definition}
\theoremstyle{remark}
\newtheorem{remark}[theorem]{Remark}
\newtheorem{example}[theorem]{Example}

\newcommand{\nrm}[1]{\left\lVert #1 \right\rVert}
\newcommand{\ML}{\operatorname{ML}}
\newcommand{\gap}{\operatorname{gap}}
\newcommand{\Z}{\mathbb{Z}}
\newcommand{\R}{\mathbb{R}}
\newcommand{\N}{\mathbb{N}}
\newcommand{\Q}{\mathbb{Q}}
\newcommand{\lean}[1]{{\def\_{\textunderscore\allowbreak}\def\.{.\allowbreak}\texttt{#1}}}
\newcommand{\LRC}{\mathrm{LRC}}

\title[A reduction of the lonely runner conjecture]{A machine-checked reduction of the lonely runner conjecture\\ to a single crux, and first-order certificates}
\author{Vico Bonfioli}
\email{vicobonfioli@gmail.com}
\date{}

\begin{document}
\emergencystretch=3em

\begin{abstract}
For positive integer speeds $v_1,\dots,v_m$ let $\ML(v)=\sup_{t}\min_i\nrm{v_it}$. The lonely
runner conjecture asserts $\ML(v)\ge 1/(m+1)$. We give a proof, checked in Lean~4 with Mathlib, that
the conjecture for all $m$ follows from three finite-range statements, each a consequence of
published or announced cases of the conjecture, together with one statement $\mathrm{FloorCruxU}$: the bound for
primitive families with at least $12$ speeds that have primitive leave-one-out subfamilies,
contain a multiple of every $j\in[2,m+1]$, carry a $\{-1,0,1\}$-relation, and admit no first-moment
certificate with a Riesz-product weight. We also verify that, given the three finite statements,
$\mathrm{FloorCruxU}$ is equivalent to the conjecture; the reduction localises the problem and does
not weaken it. Further verified results: a first-moment criterion for integer families with
relations, in the form introduced by Bedert; the fact that no absolutely continuous weight certifies a
tight family; a one-sided perturbation lemma at the times $a/(m+1)$; and the bound
$\ML\ge (m+1)/(mj+m+1)$ for the families $([1,m]\setminus\{j\})\cup\{j(m+1)\}$,
$(m+1)/2<j\le m$, $\gcd(j,m+1)=1$. For the one-entry families $([1,m]\setminus\{r\})\cup\{w\}$ we
verify two sufficient conditions for $\ML\ge 2/(2m+1)$, one at the time $a/(2m+1)$ with
$ar\equiv1$ and one at the time $1/r$; the second is the argument of Goddyn and Wong. We state
conjectures on primitive families with a speed
divisible by $m+1$, with the exact computational ground, and record examples showing that the
denominator of a certifying time is not bounded by any affine function of $m$ tested.
\end{abstract}

\maketitle

\section{Introduction}

\subsection{Conventions}
Throughout, $\nrm{x}$ is the distance from $x\in\R$ to the nearest integer. A \emph{family} is a
tuple $v=(v_1,\dots,v_m)$ of integers; $m$ is the number of \emph{speeds}, and the corresponding
runner problem has $m+1$ runners. We put
\[
  \gap(v,t)=\min_{1\le i\le m}\nrm{v_it},\qquad \ML(v)=\sup_{t\in\R}\gap(v,t).
\]
The \emph{target} for $m$ speeds is $1/(m+1)$. For $m\ge 1$ let $\LRC(m)$ denote the statement
\[
  \ML(v)\ge \frac1{m+1}\quad\text{for every } v\in\Z_{>0}^m .
\]
Repeated speeds are allowed. This is equivalent to the usual formulation with distinct speeds,
since a family with $m'<m$ distinct values has the larger target $1/(m'+1)$. A family is
\emph{tight} if $\ML(v)=1/(m+1)$.

In the Lean development a family with $m$ speeds is a function $\mathtt{Fin}\,(k+1)\to\Z$ with
$k=m-1$, and the target is written $1/(k+2)$. We translate every Lean statement to the index $m$;
the correspondence is recorded in the file \texttt{LEAN\_MAP.md} accompanying this paper, together
with the output of \verb|#print axioms| for every declaration cited. Every result labelled
\emph{Lean} below builds with no \texttt{sorry} and depends only on the axioms \texttt{propext},
\texttt{Classical.choice} and \texttt{Quot.sound}. Statements without a Lean name are either cited
from the literature or explicitly labelled as conjectures, computations, or unformalized remarks.

\subsection{Results}
Section~\ref{sec:reduction} proves (Theorem~\ref{thm:reduction}) that $\LRC(m)$ for all $m\ge3$
follows from three finite-range atoms $\mathrm{A1}$, $\mathrm{A2}$, $\mathrm{A3}$ and one crux
$\mathrm{FloorCruxU}$, and (Theorem~\ref{thm:equiv}) that under $\mathrm{A1}$--$\mathrm{A3}$ the
crux is equivalent to $\LRC(m)$ for all $m\ge 3$. The atoms concern at most $11$ speeds; they are
consequences of published or announced cases of the conjecture \cite{cp,bhk,bs,ros8,ros9,trak,st}
but are not formalized here. The conjecture is known, in the literature, for at most $13$ runners
\cite{st} (at most $14$ by the announcement \cite{allikvere}); the crux restricted to $m\le 12$
speeds is therefore known, and its open part is exactly the open part of the conjecture.

Section~\ref{sec:moment} treats first-moment certificates. The soundness of an arbitrary measure
as a certificate is Bedert's Lemma~2.1 \cite{bedert}. We verify a criterion for families with
$\{-1,0,1\}$-relations (Theorem~\ref{thm:webriesz}) which uses Bedert's Riesz product with an exact
normaliser, and a barrier (Theorem~\ref{thm:barrier}): if $\ML(v)\le\delta$ the set of times with
gap at least $\delta$ is countable, so no weight absolutely continuous with respect to Lebesgue
measure certifies a tight family. The barrier is elementary and we regard it as folklore. A related device appears in the proof
of \cite[Thm.~2.1]{beckeverett}, where the absence of a time with gap strictly above the target
makes a product of window functions vanish identically. That a family with no strict lonely time satisfies a bounded integer relation is the relation theorem of \cite{beckeverett}, building on Czerwinski \cite{czerwinski} and Giri and Kravitz \cite{gk}; no relation theorem is claimed here.

Section~\ref{sec:escape} treats certificates at a rational base point $a/q$ shifted by a common
small time $\varepsilon$. We verify a one-sided perturbation lemma (Theorem~\ref{thm:lemmaP}) and
the bound $\ML(F(m,j))\ge (m+1)/(mj+m+1)$ for the families
$F(m,j)=([1,m]\setminus\{j\})\cup\{j(m+1)\}$ (Theorem~\ref{thm:famF}). For $j=m$ the value is
Kravitz's exact value \cite[Thm.~5.1]{kravitz}; for $j<m$ the strict inequality
$\ML(F(m,j))>1/(m+1)$ follows from Zhang's classification \cite[Thm.~1.2]{zhang} combined with the
Goddyn--Wong condition \cite{gw}. What is added is the explicit bound, its witness, and the
formalization.

Section~\ref{sec:conj} states Conjectures~\ref{conj:T}--\ref{conj:T2} with their computational
ground. Section~\ref{sec:oneentry} proves (T2) for one-entry modifications of $[1,m]$ under
explicit residue conditions (Theorems~\ref{thm:oneA} and \ref{thm:oneB}), which cover all but
$42$ of the $15\,307$ primitive such families with $3\le m\le40$, $w\le 20m$ and $(m+1)\mid w$;
the remaining $42$ satisfy (T2) by exact computation. Section~\ref{sec:obstruction} records, briefly, why the certificates of
Sections~\ref{sec:moment} and \ref{sec:escape} do not by themselves reach the crux.

\section{Classes of families}\label{sec:classes}

\begin{definition}\label{def:classes}
Let $v\in\Z_{>0}^m$.
\begin{enumerate}
\item $v$ is \emph{primitive} if $\gcd(v_1,\dots,v_m)=1$.
\item $v$ is \emph{leave-one-out primitive} if for every $l$ the family $(v_i)_{i\ne l}$ is
primitive.
\item $v$ is \emph{divisor-complete} if every integer $j$ with $2\le j\le m+1$ divides some $v_i$.
\item The \emph{web} of $v$ is
$W(v)=\{\sigma\in\{-1,0,1\}^m\setminus\{0\}:\ \sum_i\sigma_iv_i=0\}$. The family is
\emph{strongly free} if $W(v)=\emptyset$ (Lean: \lean{StrongFree}).
\item $v$ is \emph{two-free} if for every $i$ the family obtained by replacing $v_i$ by $2v_i$ is
strongly free (Lean: \lean{TwoFree}).
\end{enumerate}
\end{definition}

\begin{definition}[Riesz-product union bounds]\label{def:UD}
For $v\in\Z^m$ and $\delta\in\R$ let $B_i(\delta)=\{t\in\R:\nrm{v_it}<\delta\}$ (Lean:
\lean{badSet}) and $P_v(t)=\prod_{i=1}^m(1-\cos 2\pi v_it)$. For $D\in\N$ we say that
\emph{$U_D$ fires on $v$} (Lean: \lean{UDFires}) if
\[
  \sum_{i=1}^m\int_{[0,1]\cap B_i(1/(m+1))}P_v(t)^D\,dt\ <\ \int_{[0,1]}P_v(t)^D\,dt .
\]
\end{definition}

For $D=0$ this is the Lebesgue union bound.

\section{The reduction}\label{sec:reduction}

\subsection{The atoms and the crux}
The following four statements are Lean \texttt{Prop}s; none is proved in the development.

\begin{definition}[Atoms]\label{def:atoms}
\begin{itemize}
\item[A1] (Lean: \lean{StrongFreeWindowTiny}.) For $4\le m\le 7$, every strongly free
$v\in\Z_{>0}^m$ satisfies $\ML(v)\ge1/(m+1)$.
\item[A2] (Lean: \lean{DissociatedOctetWitness}.) For every strongly free $v\in\Z_{>0}^8$ that is
not two-free there are integers $p\ge0$, $q\ge1$ such that for all $i$, with $r_i$ the least
nonnegative residue of $v_ip$ modulo $q$,
\[
  q\ \le\ 9\min(r_i,\ q-r_i),
\]
that is, $\nrm{v_ip/q}\ge 1/9$.
\item[A3] (Lean: hypothesis \lean{hwinB}.) For $1\le m\le 11$, every $v\in\Z_{>0}^m$ that is not
strongly free satisfies $\ML(v)\ge1/(m+1)$.
\end{itemize}
\end{definition}

\begin{definition}[The crux; Lean: \lean{RoomA4Compress\.FloorCruxU}]\label{def:crux}
$\mathrm{FloorCruxU}$ is the statement: for every $m\ge12$, if $\LRC(m')$ holds for all
$1\le m'<m$, then every $v\in\Z_{>0}^m$ that is primitive, leave-one-out primitive,
divisor-complete and not strongly free, and on which $U_D$ fires for no $D\in\N$, satisfies
$\ML(v)\ge1/(m+1)$.
\end{definition}

\begin{theorem}[Lean: \lean{RoomA4Compress\.universalFloor\_of\_four\_atoms\_U}]\label{thm:reduction}
$\mathrm{A1}\wedge\mathrm{A2}\wedge\mathrm{A3}\wedge\mathrm{FloorCruxU}$ implies $\LRC(m)$ for
every $m\ge3$.
\end{theorem}

The proof is a strong induction on $m$. Strongly free families are handled by a ladder argument
(Lean: \lean{separationSF\_of\_tiny\_and\_residual}, \lean{ladderSF\_le}) which uses A1 and A2;
families that are not strongly free and have $m\le11$ are A3. For $m\ge 12$ the induction step
applies, in order, Lemma~\ref{lem:ros5}, Lemma~\ref{lem:ros4}, Proposition~\ref{prop:UD}, and
finally the crux. Scale invariance reduces to primitive families.

\begin{theorem}[Lean: \lean{RoomFCUErdos\.floorCruxU\_iff\_universalFloor}]\label{thm:equiv}
Assume $\mathrm{A1}$, $\mathrm{A2}$ and $\mathrm{A3}$. Then $\mathrm{FloorCruxU}$ holds if and only
if $\LRC(m)$ holds for every $m\ge 3$.
\end{theorem}

The reverse implication is immediate, since the conclusion of the crux is an instance of
$\LRC(m)$. Theorem~\ref{thm:equiv} says that the reduction loses nothing and gains nothing: the
crux carries the whole difficulty of the conjecture at $m\ge12$ speeds.

\begin{proposition}[Lean: \lean{strongFreeWindowTiny\_of\_456}]\label{prop:A1}
A1 follows from its restriction to $5\le m\le7$. The case $m=4$ is the verified five-runner
theorem.
\end{proposition}

\begin{remark}[Status of the atoms]\label{rem:atoms}
None of A1 ($m=5,6,7$), A2, A3 is formalized. Each follows from known cases of the conjecture.
A1 and A3 are instances of $\LRC(m)$ for $m\le 7$ and $m\le 11$ respectively; these are due to
Cusick and Pomerance \cite{cp} ($m=4$), Bohman, Holzman and Kleitman \cite{bhk} ($m=5$), Barajas
and Serra \cite{bs} ($m=6$), Rosenfeld \cite{ros8} ($m=7$), Rosenfeld \cite{ros9} and
Trakulthongchai \cite{trak} ($m=8,9$), and Sungkawichai and Trakulthongchai \cite{st}
($m=10,11$; preprint). A2 follows from $\LRC(8)$ together with the fact that for rational
$\delta$ the set $\{t:\gap(v,t)\ge\delta\}$ is a finite union of closed intervals per period with
rational endpoints, so that $\ML(v)\ge\delta$ is witnessed by a rational time; this last step is
standard and is not formalized. The case $\ML(v)>\delta$ is Lean
\lean{RoomFCUCompress\.rational\_point\_of\_lt}. Consequently Theorem~\ref{thm:reduction} is a
reduction of the conjecture to the crux modulo the cited literature, and a Lean theorem only
with A1--A3 as hypotheses.
\end{remark}

\subsection{Uniform lemmas used in the step}
The following two lemmas are taken from Rosenfeld \cite{ros8}. Lemma~\ref{lem:ros5} is his
Lemma~5. His Lemma~4 states that a family without the lonely runner property has
$\operatorname{lcm}(2,\dots,m+1)\mid\prod_iv_i$; Lemma~\ref{lem:ros4} is the time-$1/j$ argument
behind it, in the form in which $j$ divides a single speed.

\begin{lemma}[Lean: \lean{RoomA4Rosenfeld\.witness\_of\_missing\_divisor}, \lean{RoomFCUErdos\.floor\_of\_not\_D}]\label{lem:ros4}
Let $v\in\Z_{>0}^m$ and $2\le j\le m+1$ with $j\nmid v_i$ for all $i$. Then
$q=j$, $p=1$ satisfy $q\le(m+1)\min(r_i,q-r_i)$ for all $i$, where $r_i\equiv v_i\pmod j$,
$0\le r_i<j$. In particular, if $m+1$ divides no $v_i$ then $\ML(v)\ge1/(m+1)$.
\end{lemma}

\begin{lemma}[Lean: \lean{RoomA4Rosenfeld\.floor\_of\_leaveOneOut\_imprimitive}]\label{lem:ros5}
Let $m\ge 3$, assume $\LRC(m-1)$, and let $v\in\Z_{>0}^m$ be primitive with
$\gcd\big((v_i)_{i\ne l}\big)\ne1$ for some $l$. Then $\ML(v)\ge1/(m+1)$.
\end{lemma}

Let $\mathcal D$ be the set of families $v\in\Z_{>0}^m$, $m\ge 3$, in which $m+1$ divides some
speed.

\begin{proposition}[Lean: \lean{RoomFCUErdos\.universalFloor\_iff\_divisible}]\label{prop:D}
$\LRC(m)$ for all $m\ge3$ is equivalent to the bound $\ML(v)\ge1/(m+1)$ for all $v\in\mathcal D$.
\end{proposition}

\begin{proposition}[Lean: \lean{RoomFCUModerator\.strictFloorD\_false}]\label{prop:strictD}
The strict bound $\ML(v)>1/(m+1)$ fails on $\mathcal D$: $v=(4,8,12)$ has $\ML(v)=1/4$.
\end{proposition}

Proposition~\ref{prop:strictD} is why Conjectures~\ref{conj:T} and \ref{conj:T2} below carry the
hypothesis $\gcd=1$.

\begin{remark}
The crux is the weakest of several forms of the last atom recorded in the development. The
earlier forms \lean{BridgeDisjunctionQPrim} and \lean{RoomA4Rosenfeld\.BridgeDisjunctionQPrimInd2}
ask, on successively smaller classes, for either the criterion of Theorem~\ref{thm:webriesz} or
a rational witness; each implies $\mathrm{FloorCruxU}$ (Lean:
\lean{RoomA4Rosenfeld\.bridgeDisjunctionQPrimInd\_of\_QPrim}, \lean{RoomA4Rosenfeld\.ind\_of\_ind2},
\lean{RoomA4Compress\.floorCruxU\_of\_ind2}). Theorem~\ref{thm:equiv} shows that no further
weakening of this type can make the crux easier than the conjecture.
\end{remark}

\section{First-moment certificates}\label{sec:moment}

\begin{proposition}[Lean: \lean{RoomA4Compress\.floor\_of\_measure\_certificate}]\label{prop:measure}
Let $v\in\Z^m$, $\delta\in\R$ and let $\nu$ be any measure on $\R$. If
$\sum_i\nu(B_i(\delta))<\nu(\R)$, then $\delta\le\ML(v)$.
\end{proposition}

This is Bedert's Lemma~2.1 \cite{bedert}, stated for an arbitrary measure. A Dirac mass recovers
rational witnesses, counting measure on a coset $t_0+q^{-1}\Z$ recovers Lemma~\ref{lem:ros5}-type
arguments, and a weight recovers the following.

\begin{proposition}[Lean: \lean{RoomA4Compress\.floor\_of\_integral\_unionBound}, \lean{RoomA4Compress\.floor\_of\_UD}]\label{prop:UD}
Let $v\in\Z^m$, $\delta\in\R$, and $F:\R\to\R$ continuous and nonnegative. If
$\sum_i\int_{[0,1]\cap B_i(\delta)}F<\int_{[0,1]}F$ then $\delta\le\ML(v)$. In particular, if $U_D$
fires on $v$ for some $D\in\N$ then $\ML(v)\ge1/(m+1)$.
\end{proposition}

\subsection{The web criterion}
For $v\in\Z^m$ define the rationals
\[
  M(v)=1+\sum_{\sigma\in W(v)}\Big(-\frac12\Big)^{|\operatorname{supp}\sigma|},\qquad
  M_l(v)=1+\sum_{\substack{\sigma\in W(v)\\ \sigma_l=0}}\Big(-\frac12\Big)^{|\operatorname{supp}\sigma|}
\]
(Lean: \lean{webMassQ}, \lean{webMassDelQ}). $M(v)$ is the constant Fourier coefficient of $P_v$,
and $M_l(v)$ that of $P_{(v_i)_{i\ne l}}$.

\begin{theorem}[Lean: \lean{floor\_of\_rieszSharpFiresQ}]\label{thm:webriesz}
Let $v\in\Z^m$. If $M(v)>0$ and
\[
  99225\Big(\sum_{l=1}^m\big(2M_l(v)-M(v)\big)+M(v)\Big)\ \le\ 10000\,(m+1)^2\,M(v),
\]
then $\ML(v)\ge1/(m+1)$.
\end{theorem}

The hypothesis is the predicate \lean{rieszSharpFiresQ}; the constant $99225/10000$ exceeds
$\pi^2$. The proof applies Proposition~\ref{prop:measure} to the probability measure with density
$P_v/M(v)$ and bounds the mass of each $B_l$ through the subfamily coefficient $M_l(v)$. The
method, including the expansion of $P_v$ over $\{-1,0,1\}^m$ with weights
$(-1/2)^{|\operatorname{supp}\sigma|}$, is Bedert's \cite[Lemma~4.1]{bedert}, where $v$ is
dissociated and $M(v)=1$. The additions are the normaliser $M(v)$ for families with relations and
the exact per-coordinate term in place of a supremum bound.

\begin{corollary}[Lean: \lean{rieszSharpFiresQ\_of\_subfamilySum\_le}]\label{cor:web}
Let $m\ge9$ and $v\in\Z^m$ with $M(v)>0$ and $\sum_lM_l(v)\le m\,M(v)$. Then the hypothesis of
Theorem~\ref{thm:webriesz} holds, and $\ML(v)\ge1/(m+1)$.
\end{corollary}

\begin{remark}
If $v$ is strongly free then $W(v)=\emptyset$, so $M(v)=M_l(v)=1$ and Corollary~\ref{cor:web}
gives $\ML(v)\ge 1/(m+1)$ for strongly free families with $m\ge 9$. The identification
$W(v)=\emptyset$ is immediate from the definitions and is not a separate Lean declaration.
Bedert's supremum bound, read at the target $1/(m+1)$, gives the same conclusion from $m=18$.
\end{remark}

\subsection{The barrier}
For $\delta\in\R$ put $G_\delta(v)=\{t\in\R:\ \nrm{v_it}\ge\delta\text{ for all }i\}$.

\begin{theorem}\label{thm:barrier}
Let $v\in\Z^m$ with all $v_i\ne0$.
\begin{enumerate}
\item (Lean: \lean{RoomA4Barrier\.goodSet\_countable}.) If $\ML(v)\le\delta$ then $G_\delta(v)$ is
countable.
\item (Lean: \lean{RoomA4Barrier\.integral\_le\_sum\_badSet}.) If $\ML(v)\le\delta$ and $F\ge0$ is
integrable on $[0,1]$, then $\int_{[0,1]}F\le\sum_i\int_{[0,1]\cap B_i(\delta)}F$.
\item (Lean: \lean{RoomA4Barrier\.lt\_ML\_of\_UDFires}.) If $U_D$ fires on $v$ then
$\ML(v)>1/(m+1)$.
\item (Lean: \lean{RoomA4Barrier\.not\_UDFires\_of\_tight}.) If $v$ is tight, $U_D$ fires on $v$
for no $D$.
\item (Lean: \lean{RoomA4Barrier\.not\_UDFires\_consec}.) $U_D$ fires on $(1,2,\dots,m)$ for no $D$.
\end{enumerate}
\end{theorem}

Thus every tight family satisfies the last hypothesis of the crux, and Proposition~\ref{prop:UD}
only ever certifies the strict inequality. Whether a tight family lies in the class of
Definition~\ref{def:crux} is open; Conjecture~\ref{conj:T} would exclude it.

\section{Escape certificates}\label{sec:escape}

\begin{proposition}[Lean: \lean{RoomFCUModerator\.ML\_ge\_of\_escape}]\label{prop:escape}
Let $v\in\Z^m$, $q,a\in\Z$ with $q>0$, and $\varepsilon,\delta\in\R$ with $\delta>0$. Suppose that
for every $i$ there are $N_i,r_i\in\Z$ with $v_ia=qN_i+r_i$ and
\[
  \delta\ \le\ \frac{r_i}{q}+v_i\varepsilon\ \le\ 1-\delta .
\]
Then $\delta\le\ML(v)$.
\end{proposition}

The time used is $a/q+\varepsilon$. The statement is elementary and the pattern is standard
\cite{kravitz}. With $\varepsilon=0$ and $\delta=A/B$ it takes a purely arithmetic form (Lean:
\lean{RoomFCUCompress\.ML\_ge\_of\_residues}). Its converse also holds with $q=1$ and
$\varepsilon$ a maximiser (Lean: \lean{RoomFCUCompress\.lonely\_iff\_escape}); the proposition
therefore has content only together with a bound on $q$, which is the subject of
Section~\ref{sec:obstruction}.

\begin{theorem}[Lemma P; Lean: \lean{RoomFCUCompress\.lemmaP\_quant}, \lean{RoomFCULemmaP\.lemmaP}]\label{thm:lemmaP}
Let $m\ge3$, $q=m+1$ and $v\in\Z_{>0}^m$. Suppose there is an index $i_0$ such that
\begin{enumerate}
\item $q\mid v_{i_0}$;
\item for every $i\ne i_0$, $q\nmid v_i$ and $v_i<v_{i_0}$;
\item there is $u$ with $\gcd(u,q)=1$ and $v_i\not\equiv u\pmod q$ for all $i$.
\end{enumerate}
Put $W=v_{i_0}$. Then
\[
  \ML(v)\ \ge\ \frac{2W}{q(2W-1)}\ >\ \frac1{m+1}.
\]
\end{theorem}

The witness is $t=a/q-2/(q(2W-1))$ with $au\equiv1\pmod q$: the largest speed moves from $0$ to
distance $2W/(q(2W-1))$, while every other speed has residue $r\in[2,q-2]$ at $a/q$ and moves by
less than the available margin. We found no statement of this lemma in the literature consulted
(\cite{kravitz,zhang,ps}).

\begin{theorem}[Lean: \lean{RoomFCUCompress\.famF\_via\_escape}, \lean{RoomFCURamanujan\.famF\_lonely}]\label{thm:famF}
Let $m,j$ be positive integers with $(m+1)/2<j\le m$ and $\gcd(j,m+1)=1$, and put
$F(m,j)=([1,m]\setminus\{j\})\cup\{j(m+1)\}$ and $D=mj+m+1$. There is an integer $s\ge0$ such that
the time $t=(sD+1)/(jD)$ satisfies $\nrm{xt}\ge(m+1)/D$ for every $x\in F(m,j)$. Consequently every
family $v$ all of whose speeds lie in $F(m,j)$ satisfies
\[
  \ML(v)\ \ge\ \frac{m+1}{mj+m+1}\ >\ \frac1{m+1}.
\]
\end{theorem}

For $j=m$, $F(m,m)=\{1,\dots,m-1,m(m+1)\}$ and the bound is Kravitz's value
$\ML(1,\dots,n-1,ns)=s/(ns+1)$ at $n=m$, $s=m+1$, where equality holds \cite[Thm.~5.1]{kravitz}.
For $j<m$, $F(m,j)$ is $[n-1]$ with $n=m+1$ and the speed $r=j$ replaced by $\mu r$,
$\mu=m+1$. Goddyn and Wong \cite[Thm.~2.3]{gw} show that such a family is tight only if
$\gcd(r,b)>1$ for every $b$ with $n-r\le b<\mu(n-r)$; here $b=n-r=m+1-j$ gives
$\gcd(j,m+1-j)=\gcd(j,m+1)=1$, so $\ML(F(m,j))>1/(m+1)$ was known (see also Zhang's
classification \cite[Thm.~1.2]{zhang}). Theorem~\ref{thm:famF} makes this
quantitative. The value has the shape $s/(ns+k)$ with $(n,s,k)=(j,m+1,m+1-j)$, the shape of the
amended spectrum conjecture of Fan and Sun \cite{fansun}; we have not checked that their
parameter ranges cover these values.

\begin{conjecture}\label{conj:Feq}
Under the hypotheses of Theorem~\ref{thm:famF}, $\ML(F(m,j))=(m+1)/(mj+m+1)$.
\end{conjecture}

Ground: exact rational computation of $\ML(F(m,j))$ agrees in all $499$ admissible pairs with
$m\le56$ (Appendix~\ref{app:data}).

\begin{proposition}[Lean: \lean{RoomFCUGoddyn\.fast\_runner\_floor}, \lean{RoomFCUGoddyn\.cruxU\_fast}]\label{prop:fast}
Let $m\ge 2$, $v\in\Z_{>0}^m$ and $l$ an index such that $mv_i\le v_l$ for all $i\ne l$. If
$\ML\big((v_i)_{i\ne l}\big)\ge1/m$ then $\ML(v)\ge1/(m+1)$.
\end{proposition}

This is the case $L=1/m$, $\varepsilon=1/(m(m+1))$ of Kravitz's Proposition~8.1 \cite{kravitz}. It gives the crux for every family with one
speed at least $m$ times all others, for example $F(m,m)$, and it is not an instance of
Proposition~\ref{prop:escape} with bounded $q$.

\section{Conjectures on the divisible class}\label{sec:conj}

All computations in this section use exact integer arithmetic and the standard fact that
$\ML(v)$ is attained at a time $c/(v_i+v_j)$ with $i\le j$ and $c\in\Z$; this fact is not
formalized. A \emph{box} $(m,H)$ is the set of primitive $m$-element subsets of $[1,H]$, enumerated
completely.

\begin{conjecture}[(T)]\label{conj:T}
If $v\in\Z_{>0}^m$ is primitive and tight, then $m+1$ divides no $v_i$.
\end{conjecture}

\begin{conjecture}[(T$'$)]\label{conj:Tprime}
If $v\in\Z_{>0}^m$ is primitive and tight, then every $t$ with $\gap(v,t)=1/(m+1)$ lies in
$\frac1{m+1}\Z$.
\end{conjecture}

(T$'$) implies (T): at $t=a/(m+1)$ a speed divisible by $m+1$ is at distance $0$. For
$v=(1,\dots,m)$, (T$'$) is contained in Kravitz's Proposition~10.4 \cite{kravitz}: the equality
times are the $a/(m+1)$ with $\gcd(a,m+1)=1$. By Lemma~\ref{lem:ros4},
(T) is equivalent to the statement that no primitive tight family is divisor-complete with
$m+1$ in the divisor set; in particular (T) would exclude tight families from the class of the
crux. The gcd hypothesis cannot be dropped: $5\cdot(1,2,3,4)$ is tight and every speed is
divisible by $5$.

\begin{proposition}[Lean: \lean{RoomFCUHilbert\.gw\_not\_dvd}, \lean{RoomFCUHilbert\.gw\_family\_not\_dvd}]\label{prop:gw}
Let $0<r<n$ and $\mu\ge1$ be integers such that $\gcd(r,b)>1$ for every integer $b$ with
$n-r\le b<\mu(n-r)$. Then $n\nmid r\mu$.
\end{proposition}

The hypothesis is the Goddyn--Wong condition for accelerating the runner $r$ of $[n-1]$ to $r\mu$
\cite[Thm.~2.3]{gw}, and the argument is theirs: it is the last step of the proof of
\cite[Thm.~3.1]{gw}, which shows $n\nmid m_rr$ for every accelerated speed $m_rr$ in order to use
the time $1/n$. We have checked every tight instance listed or constructed in \cite{gw} against
(T). The sporadic list there (all tight vectors with $n\le 20$ and maximum speed at most $40$:
$\{1,3,4,7\}$, $\{1,3,4,5,9\}$, $\{1,4,5,6,7,11,13\}$, $\{1,2,3,4,5,7,12\}$, $[13]_{12\mapsto24}$,
$[19]_{18\mapsto36}$) has no speed divisible by $n$; the families of \cite[Thms.~2.3, 3.1]{gw}
have none by the argument just quoted; the exceptional case $(n,r,\mu)=(3,1,4)$ of
\cite[Thm.~2.3]{gw} is $(4,2)$, which is not primitive and has no speed divisible by $3$. The
two further sporadic tight one-entry families of Zhang \cite[Thm.~1.2]{zhang}, $\{1,3,4,7\}$ and
$\{1,3,4,5,9\}$, are the first two of this list. Hence no published tight family contradicts (T).
The deduction from Proposition~\ref{prop:gw} to those instances rests on the cited theorems and is
not formalized.

\begin{conjecture}[T2]\label{conj:T2}
If $v\in\Z_{>0}^m$, $m\ge 3$, is primitive and $m+1$ divides some $v_i$, then
$\ML(v)\ge 2/(2m+1)$.
\end{conjecture}

\begin{remark}\label{rem:T2}
(T2) implies $\LRC(m)$: for $v\in\Z_{>0}^m$ with $g=\gcd(v)$, $\ML(v)=\ML(v/g)$; if $m+1$ divides
some $v_i/g$ apply (T2), otherwise Lemma~\ref{lem:ros4}. This argument is short and is not
formalized. (T2) fails without $\gcd=1$ (Proposition~\ref{prop:strictD}) and fails for primitive
families outside $\mathcal D$: $\ML(1,5,6,11,16,17)=5/33<2/13$ (computed).
\end{remark}

\subsection{Computational ground}
Table~\ref{tab:boxes} lists the complete enumerations. In every box, every primitive family with
$\ML(v)<1/m$ was retained and its value computed exactly. Since $2/(2m+1)<1/m$, a counterexample
to (T2) in the box would have been retained.

\begin{table}[h]
\centering
\begin{tabular}{rrrl}
\toprule
$m$ & $H$ & primitive families & tight families found\\
\midrule
3 & 1500 & 467\,278\,456 & $[1,3]$\\
4 & 500 & 2\,380\,658\,191 & $[1,4]$, $\{1,3,4,7\}$\\
5 & 200 & 2\,450\,913\,356 & $[1,5]$, $\{1,3,4,5,9\}$\\
6 & 110 & 2\,110\,853\,351 & $[1,6]$\\
7 & 60 & 384\,092\,920 & $[1,7]$, $\{1,2,3,4,5,7,12\}$, $\{1,4,5,6,7,11,13\}$\\
8 & 48 & 376\,600\,645 & $[1,8]$\\
9 & 40 & 273\,270\,205 & $[1,9]$\\
10 & 36 & 254\,143\,032 & $[1,10]$\\
\bottomrule
\end{tabular}
\caption{Complete enumerations of primitive $m$-subsets of $[1,H]$.}\label{tab:boxes}
\end{table}

In all boxes of Table~\ref{tab:boxes}: no family has $\ML(v)<1/(m+1)$; no tight family has a speed
divisible by $m+1$; every maximiser of every tight family has denominator $m+1$. This is the
ground for (T) and (T$'$). An independent enumeration with a different program (boxes
$(4,400)$, $(5,200)$, $(6,100)$, $(7,64)$, $(8,50)$, $(9,42)$, $(10,38)$) found the same tight
families.

For (T2) the minimum of $\ML$ over the primitive members of $\mathcal D$ in the box is exactly
$2/(2m+1)$ for $m=3,\dots,7$, and is $4/33$, $4/37$, $3/31$ for $m=8,9,10$, all above
$2/(2m+1)$. An earlier search at $m=12$ over heights at most $200$ reported $\ML=2/25$ for
$(1,4,6,7,10,11,12,13,16,17,19,23)$. That value is wrong and the claim that the bound is attained at $m=12$ is withdrawn:
the program that produced it enumerated the times $c/(v_i+v_j)$ with $c$ coprime to $v_i+v_j$ only, and so missed the maximiser
$t=1/9$, whose denominator is not a sum of two speeds. The exact value is $\ML=1/9>2/25$ (exact enumeration over all numerators,
\texttt{code/obstructions/mlcheck}; a scan of all $t=p/q$ with $q<300$ in exact rational arithmetic gives the same value). Equality $\ML=2/(2m+1)$
in $\mathcal D$ is therefore recorded only for $m=3,\dots,7$ (Table~\ref{tab:boxes}).

For $m=14,16,18,20$ simulated annealing over primitive members of $\mathcal D$ (not complete) ran
with the parameters of Table~\ref{tab:t2}. No value below $2/(2m+1)$ was found. At height $5000$
the minimum found was $\ML(F(m,m))=(m+1)/(m^2+m+1)$ in every case.

\begin{table}[h]
\centering
\begin{tabular}{rrrrl}
\toprule
$m$ & max height & time (s) & evaluations & least $(2m+1)\ML$ found\\
\midrule
14 & 5000 & 720 & 70\,353\,611 & $2.0616$ \ ($15/211$)\\
16 & 5000 & 720 & 25\,184\,829 & $2.0549$ \ ($17/273$)\\
18 & 5000 & 720 & 15\,945\,569 & $2.0496$ \ ($19/343$)\\
20 & 5000 & 720 & 13\,549\,085 & $2.0451$ \ ($21/421$)\\
14 & 60 & 180 & 63\,550\,117 & $2.0714$ \ ($1/14$)\\
16 & 60 & 180 & 48\,313\,493 & $2.0625$ \ ($1/16$)\\
18 & 60 & 180 & 37\,979\,054 & $2.0556$ \ ($1/18$)\\
20 & 60 & 180 & 39\,240\,525 & $2.0500$ \ ($1/20$)\\
\bottomrule
\end{tabular}
\caption{Searches for counterexamples to (T2); one run per row, $4$ threads, fixed seeds.}\label{tab:t2}
\end{table}

Tight instances in the boxes are few, so the independent evidence for (T) and (T$'$) is thin.
The value $(2m+1)(m+1)/(m^2+m+1)$ tends to $2$, so (T2), if true, is asymptotically sharp along
$F(m,m)$.

\section{(T2) on one-entry families}\label{sec:oneentry}

For $1\le r\le m$ and $w\in\Z$ let
\[
  V(m,r,w)=\big([1,m]\setminus\{r\}\big)\cup\{w\},
\]
the family $(1,\dots,m)$ with the entry $r$ replaced by $w$ (Lean: \lean{T2Sub\.oneEntry}). No
hypothesis on $w$ other than those stated is needed below; in particular $w$ may coincide with
an entry, and $(m+1)\mid w$ is not assumed.

\begin{theorem}[Lean: \lean{T2Sub\.oneEntry\_ML\_ge}]\label{thm:oneA}
Let $1\le r\le m$, and let $q,a$ be integers with $q>m$, $q\mid ar-1$, and $q>m+r$ or $q=2r$. If
$w\not\equiv 0,\ r,\ -r\pmod q$, then $\ML(V(m,r,w))\ge 2/q$.
\end{theorem}

\begin{proof}
Take $t=a/q$. Since $a$ is inverse to $r$ modulo $q$, $xa\equiv 0,\pm1\pmod q$ holds exactly when
$x\equiv 0,\pm r\pmod q$. For $x\in[1,m]\setminus\{r\}$: $0<x<q$; $x\equiv r$ forces $x=r$
because $|x-r|<q$; and $x\equiv -r$ forces $x+r=q$, impossible if $q>m+r$ and giving $x=r$ if
$q=2r$. For $x=w$ the three congruences are excluded by hypothesis. So every $xa$ has residue in
$[2,q-2]$ and $\nrm{xt}\ge 2/q$. The Lean proof goes through Lean
\lean{RoomFCUCompress\.ML\_ge\_of\_residues}.
\end{proof}

\begin{corollary}[Lean: \lean{T2Sub\.oneEntry\_T2}]\label{cor:oneA}
If $1\le r\le m$, $\gcd(r,2m+1)=1$ and $w\not\equiv 0,\pm r\pmod{2m+1}$, then
$\ML(V(m,r,w))\ge 2/(2m+1)$.
\end{corollary}

\begin{theorem}[Lean: \lean{T2Sub\.oneEntry\_ML\_ge\_inv}, \lean{T2Sub\.oneEntry\_T2\_inv}]\label{thm:oneB}
If $1\le r\le m<2r$ and $r\nmid w$, then $\ML(V(m,r,w))\ge 1/r\ge 2/(2m+1)$.
\end{theorem}

\begin{proof}
At $t=1/r$ every speed $x$ has $\nrm{xt}\ge1/r$ unless $r\mid x$; the only multiple of $r$ in
$[1,m]$ is $r$, which has been removed, and $r\nmid w$. Finally $2r\le 2m<2m+1$.
\end{proof}

Theorem~\ref{thm:oneB} is not new. It is the first case of the proof of Goddyn and Wong's
Lemma~2.2 \cite{gw} (a proper subfamily $[n-1]\setminus\{x\}$ with $2x\ge n$ has
$\ML\ge1/x$, by the time $1/x$), and the same time is used in the proof of Zhang's
Lemma~2.3 \cite{zhang} (if $2r>n-1$ and $[n-1]_{r\to w}$ is tight then $r\mid w$). We include it
because it covers residue classes of $w$ that Theorem~\ref{thm:oneA} misses.
Theorem~\ref{thm:oneA} uses the standard base point $a/q$ with $q$ just above $m$; we did not
find it stated in \cite{gw,zhang,kravitz}, but it is a base-point argument of the type of
Lemma~\ref{lem:ros4} and Proposition~\ref{prop:escape}, and we do not claim it as a new method. Zhang
\cite[Thm.~2.9]{zhang} computes $\ML([n-1]\setminus\{2\})=2/(2\lfloor n/2\rfloor+3)$ exactly;
Theorem~\ref{thm:oneA} gives only a lower bound, for all $r$.

Combined with the classification of tight one-entry families \cite[Thm.~1.2]{zhang} and
Proposition~\ref{prop:gw}, no one-entry family with $(m+1)\mid w$ is tight, so (T) holds on this
class (a consequence of the cited theorem, not formalized). (T2) asks for the quantitative
bound $2/(2m+1)$.

\subsection*{Computational ground}
Let $\mathcal O$ be the set of triples $(m,r,w)$ with $3\le m\le 40$, $1\le r\le m$,
$(m+1)\mid w$, $m+1\le w\le 20m$ and $\gcd V(m,r,w)=1$; $|\mathcal O|=15\,307$. For every triple
we computed $\ML(V(m,r,w))$ exactly (Appendix~\ref{app:data}). Results:
\begin{enumerate}
\item $\min_{\mathcal O}(2m+1)\ML=2$; no member violates (T2).
\item Equality $\ML=2/(2m+1)$ holds exactly for $\{1,3,4\}$ ($m=3$), $\{1,3,4,5\}$ ($m=4$) and
$\{1,3,4,5,18\}$ ($m=5$), that is, $(m,r,w)=(3,2,4),(4,2,5),(5,2,18)$.
\item Theorem~\ref{thm:oneA} with some $q\in[m+2,2m+1]$ or Theorem~\ref{thm:oneB} applies to
all but $42$ triples. These are listed in Table~\ref{tab:uncov}; for each, (T2) holds by the
computed value.
\end{enumerate}
Thus (T2) on one-entry families follows from Theorems~\ref{thm:oneA} and \ref{thm:oneB} for
every triple satisfying their residue conditions, for all $m$ and without the bound $w\le 20m$;
outside those residue classes it is supported by computation only.

\section{Obstructions}\label{sec:obstruction}

\begin{remark}[Base points]
Lemma~\ref{lem:ros4}, Theorem~\ref{thm:lemmaP} and Theorem~\ref{thm:famF} are instances of
Proposition~\ref{prop:escape}: a base point $a/q$ and a common shift $\varepsilon$, followed by a
residue check. By Proposition~\ref{prop:escape} and its converse this form is complete for every
family, so all content lies in bounding $q$ in terms of $m$. Proposition~\ref{prop:fast} is an
instance only with $q=1$, and the first-moment certificates of Section~\ref{sec:moment} are not
of this form; by Theorem~\ref{thm:barrier} the latter cannot certify a tight family.
\end{remark}

\begin{remark}[Level drift]
Let $m\ge12$ and let $\mathcal C_m$ be the class of Definition~\ref{def:crux}. Computation shows
that no affine bound $q\le 2m+c$ with $c\in\{-m+1,1,2\}$ suffices, on $\mathcal C_m$, for
the variants of Proposition~\ref{prop:escape} tested (window certificates at a single level
$q$, first-order shifts, exact scans). Examples (membership in
$\mathcal C_m$ and all levels computed exactly; not formalized):
\begin{itemize}
\item $m=12$: $(1,4,5,6,7,9,11,13,40,60,238,2731)$ has no first-order window certificate with
$q\le m+1=13$, even allowing both signs and an exact scan of the multiples of $q$; the first such
level is $22$.
\item $m=16$: $(1,5,6,7,8,9,11,13,17,19,23,25,29,31,366,1680)$ has no exact window certificate
at a level $q$ with $m+1<q\le 2m+2=34$; the first such level is $35$. This family has a
certificate at $q=2$, so it bounds only the window variant.
\item $m=18$: $(1,3,4,5,7,11,13,17,19,23,31,70,74,116,180,272,3986,4943)$ has no certificate,
first-order or exact, at $q\le 2m+1=37$; the first is a point certificate at $q=38$, and
$\ML=28/191$.
\item $m=18$: $(1,3,4,5,7,11,13,16,17,19,23,25,29,31,108,185,210,3685)$ has its first
certificate at $q=46>2m+2$.
\end{itemize}
In each case a certificate exists at a larger level, so $\ML(v)\ge1/(m+1)$. The families
consist of a small core together with two or three large speeds, and the level at which a
certificate first appears was not predicted by any of the invariants tested
(Appendix~\ref{app:data}).
\end{remark}

\begin{proposition}[Lean: \lean{RoomFCUMss\.halfShift\_not\_lonely}]\label{prop:half}
There is no $t\in\R$ with
\[
  \nrm{t},\ \nrm{2t-\tfrac12},\ \nrm{3t-\tfrac12},\ \nrm{4t-\tfrac12},\ \nrm{6t-\tfrac12},\ \nrm{9t}
  \ \ge\ \tfrac17 .
\]
\end{proposition}

With zero phases the speeds $(1,2,3,4,6,9)$ satisfy $\LRC(6)$ \cite{bs}. Hence an argument for the
conjecture must use that the starting phases are all zero, not merely that they lie in
$\frac12\Z$. Counterexamples to the conjecture with general phases are due to Blanco, Criado and
Santos \cite{bcs}, from $n=5$ in their indexing; see also \cite{poliakova} for their growth. Only the instance
with phases in $\{0,\frac12\}$ is not taken from the literature.

\appendix
\section{Computational ground and reproduction}\label{app:data}

All programs are in Rust with exact integer arithmetic and are in
\texttt{code/reduction/<name>/}; the file \texttt{code/reduction/README.md} gives the exact
command reproducing each number.
\begin{itemize}
\item Table~\ref{tab:boxes}, maximiser denominators, (T2) minima for $m\le10$: \texttt{fcuspec}.
Independent enumeration: \texttt{gdn}; the smaller boxes of a third program \texttt{ecalc} for
$m\le8$ agree.
\item Table~\ref{tab:t2}: \texttt{t2hunt}, seeds $101m+1$ (height $5000$) and $101m+2$ (height
$60$); logs record the reported minima.
\item Conjecture~\ref{conj:Feq}: exact $\ML(F(m,j))$ for all admissible $(m,j)$, $m\le 56$:
\texttt{rmj}.
\item Level-drift examples: \texttt{mdr}, an independent checker; the values of $\ML$ used in the text are recomputed by
\texttt{code/obstructions/mlcheck}, which enumerates all numerators (the function \texttt{ml} of \texttt{mdr} originally skipped numerators not coprime
to the denominator and could return a value below $\ML$; the archived source is corrected); re-checked with \texttt{modchk}. Invariant fits of
the certificate levels: \texttt{levelfit}.
\item Section~\ref{sec:oneentry}: \texttt{t2sweep} (sweep of $\mathcal O$, coverage by
Theorems~\ref{thm:oneA} and \ref{thm:oneB}, Table~\ref{tab:uncov}).
\end{itemize}


\begin{table}[h]
\centering\small
\begin{tabular}{rrrr@{\qquad}rrrr@{\qquad}rrrr}
\toprule
$m$ & $r$ & $w$ & $\ML$ & $m$ & $r$ & $w$ & $\ML$ & $m$ & $r$ & $w$ & $\ML$\\
\midrule
$3$ & $1$ & $36$ & $5/13$ & $4$ & $4$ & $40$ & $10/41$ & $7$ & $7$ & $112$ & $16/113$\\
$3$ & $2$ & $12$ & $2/5$ & $5$ & $2$ & $90$ & $10/47$ & $8$ & $6$ & $108$ & $18/113$\\
$3$ & $2$ & $16$ & $8/19$ & $5$ & $3$ & $30$ & $5/17$ & $8$ & $8$ & $144$ & $18/145$\\
$3$ & $2$ & $28$ & $14/31$ & $5$ & $4$ & $48$ & $4/17$ & $9$ & $8$ & $160$ & $20/163$\\
$3$ & $2$ & $40$ & $20/43$ & $5$ & $4$ & $84$ & $7/29$ & $9$ & $9$ & $180$ & $20/181$\\
$3$ & $2$ & $44$ & $22/47$ & $5$ & $5$ & $60$ & $12/61$ & $13$ & $12$ & $84$ & $7/89$\\
$3$ & $2$ & $56$ & $28/59$ & $6$ & $2$ & $119$ & $27/124$ & $13$ & $12$ & $168$ & $14/173$\\
$3$ & $3$ & $24$ & $8/25$ & $6$ & $6$ & $84$ & $14/85$ & $13$ & $12$ & $252$ & $21/257$\\
$3$ & $3$ & $60$ & $20/61$ & $7$ & $4$ & $56$ & $14/61$ & $19$ & $18$ & $180$ & $2/37$\\
$4$ & $1$ & $55$ & $19/59$ & $7$ & $6$ & $24$ & $4/29$ & $19$ & $18$ & $360$ & $4/73$\\
$4$ & $2$ & $65$ & $19/68$ & $7$ & $6$ & $48$ & $8/53$ & $25$ & $24$ & $312$ & $13/317$\\
$4$ & $2$ & $70$ & $20/73$ & $7$ & $6$ & $72$ & $12/77$ & $31$ & $30$ & $480$ & $16/487$\\
$4$ & $3$ & $45$ & $15/47$ & $7$ & $6$ & $96$ & $16/101$ & $32$ & $30$ & $330$ & $11/337$\\
$4$ & $3$ & $75$ & $25/77$ & $7$ & $6$ & $120$ & $4/25$ & $37$ & $36$ & $684$ & $19/689$\\
\bottomrule
\end{tabular}
\caption{The $42$ triples of $\mathcal O$ covered by neither Theorem~\ref{thm:oneA} ($q\le 2m+1$)
nor Theorem~\ref{thm:oneB}, with the exact value of $\ML(V(m,r,w))$.}\label{tab:uncov}
\end{table}

\bibliographystyle{plain}
\bibliography{refs}

\end{document}
